\section{From the integral lattice to an embedded link}\label{sec:fourfold-arithmetic}

We now compare the two actual Milnor lattices supplied by
Theorem~\ref{thm:tube}. Two boundary classes of the cone page give
rational Seifert isometries that are integral at complementary sets
of primes. We compare these two maps and use strong approximation to
replace their local choices by one integral isometry. The final
subsection then identifies the embedded links. Throughout, we use the
full integral lattice formula, including its mod-four boundary
condition and the signs of both filling forms.

The parameters \(r,n_0,n_1,a,b,\alpha,\beta\) remain the single
choice fixed in Lemma~\ref{lem:simultaneous-parameters}. Its analytic
bounds are already available. We begin with the remaining
divisibility facts and positive indices needed for the real-sign
calculation.

\begin{lemma}[Divisibility and positive indices]\label{lem:parameter-indices}
For the parameters fixed in Lemma~\ref{lem:simultaneous-parameters}, put
\[
 g=\gcd(a,b),\qquad m=4ab+a+b.
\]
Then \(a,b\equiv2\pmod5\), \(5\nmid g\), and
\(g\ge 2^{r-1}>5\). For each \(q\in\{1,4\}\), there are positive
integers \(i,j\) satisfying
\[
 i\equiv j\pmod4,\qquad \frac{bi+aj}{m}=\frac q5 .
\]
\end{lemma}

\begin{proof}
Both \(n_0,n_1\) are \(2\pmod5\), so
\((n_j^2-1)/4\equiv2\pmod5\). Thus \(a,b\equiv2\pmod5\)
and \(5\nmid g\). Also \(2^r\mid n_j-1\) and \(2\mid n_j+1\),
whence \(2^{r-1}\mid a,b\) and \(g\ge2^{r-1}>5\).

Write \(a=ga'\), \(b=gb'\), with \(\gcd(a',b')=1\).
The desired index equations, on writing \(i=j+4t\), become
\[
 (a'+b')j+4b't=\frac{qm}{5g}.
\]
Reduction modulo five gives \(m=4ab+a+b\equiv0\pmod5\);
since \(5\nmid g\), we have \(5\mid m/g\). Let
\(d=\gcd(a'+b',4b')\). The equality \(\gcd(a'+b',b')=1\)
gives \(d\mid4\), and \(d\) also divides
\(m/g=4ga'b'+a'+b'\). Hence \(d\mid m/(5g)\), so the
displayed equation has integer solutions for \(q=1,4\).
Its homogeneous solutions change \(j\) by \(4b'/d\); choose a
solution with \(1\le j\le4b'/d\le4b'\). For \(q=1\),
\[
 \frac{m}{5g}
   =\frac{4ga'b'+a'+b'}5
   >4a'b'
\]
because \(g>5\). Thus
\(i=(qm/(5g)-a'j)/b'>0\) for both values of \(q\).
\end{proof}

Put $k=4$.  For $\nu\in\{0,1\}$ let
$\pi_\nu\colon\mathcal D_\nu\to\mathcal H$ be the degree-$k$ map of
plane-curve Milnor pages constructed above.  Write
\[
 \begin{aligned}
 H&=H_1(\mathcal H;\mathbb Z),&
 D_\nu&=H_1(\mathcal D_\nu;\mathbb Z),\\
 K_\nu&=\ker(\pi_{\nu*}\colon D_\nu\to H),&
 \operatorname{pr}_\nu l
   &=l-\frac1k\pi_\nu^!\pi_{\nu*}l.
 \end{aligned}
\]
The transfer identities are $\pi_{\nu*}\pi_\nu^!=k$ and
$s_{d_\nu}(\pi_\nu^!x,l)=s_h(x,\pi_{\nu*}l)$ in both argument orders.
In particular $\operatorname{pr}_\nu$ projects rationally onto $K_\nu$
along $\pi_\nu^!H_{\mathbb Q}$, orthogonally for the Seifert forms.
The plane-curve argument has supplied an integral Seifert isometry
\begin{equation}\label{eq:arith-plane-isometry}
 \mathcal I\colon(D_0,s_{d_0})\xrightarrow{\ \sim\ }(D_1,s_{d_1}),
\end{equation}
and has shown that their monodromies are semisimple at eigenvalue
$1$ and at primitive fifth roots.  These facts are separate inputs to the
present section.

Let $Q$ be a Milnor page of $P=w(y^a+z^b)$ on the cone
$S=\{yz=w^4\}$, let $E=H_1(Q;\mathbb Z)$, and let
$\tau\colon E\to H_1(\operatorname{link}(S))=\mathbb Z/4$ be induced by
weighted radial projection $Q\to\operatorname{link}(S)$, equivalently
by page inclusion in the link open-book model. The rational form $e$
on $E$ is one fourth of the upstairs Seifert pairing of transferred
cycles. Theorem~\ref{thm:tube} identifies the Milnor lattice and its
Seifert form $s_\nu$, in one convention on both sides, as follows:
\begin{equation}\label{eq:arith-input-lattice}
 \begin{split}
 F_\nu&=(K_\nu)_{\mathbb Q}\oplus(H\otimes E)_{\mathbb Q},\\
 s_\nu&=(-s_{d_\nu}|_{K_\nu})\perp(s_h\otimes e),\\
 \Lambda_\nu&=
 \left\{(\operatorname{pr}_\nu l,z):
       l\in D_\nu,\ z\in H\otimes E,\quad
       \pi_{\nu*}l\equiv(1\otimes\tau)z\pmod{4H}\right\}.
 \end{split}
\end{equation}
The rational vector space $F_\nu$ is $\Lambda_\nu\otimes\mathbb Q$.
Every later localization uses this complete integral lattice with the
displayed residue condition; all tensor products and pairings use the
displayed identification.

\subsection{Two boundary vectors}\label{subsec:arith-boundary}
The graph index \(\nu\in\{0,1\}\) and the boundary index
\(j\in\{0,1\}\) are independent.

The quotient map $(U,V)\mapsto(U^4,V^4,UV)$ identifies $S$ with
$\mathbb C^2/\mu_4$, with action $(U,V)\mapsto(\omega U,\omega^{-1}V)$.
Upstairs the cone-page equation is
\begin{equation}\label{eq:arith-upstairs-page}
 \widetilde P(U,V)=UV(U^{4a}+V^{4b}).
\end{equation}
Orient the boundary class $b_0\in E$ of the $y$-axis as a complex
branch and the class $b_1$ of the $z$-axis oppositely to its complex
orientation.  The quotient lens space has $H_1=\mathbb Z/4$; with the
positive $y$-axis class as generator, both have $\tau(b_j)=1$.

\begin{lemma}[Boundary pairings]\label{lem:arith-boundary}
The classes $b_0,b_1$ are fixed by the page monodromy, and
\begin{equation}\label{eq:arith-boundary-pairings}
 \begin{aligned}
 e(b_j,b_j)&=-\frac{n_j^2}{4},&
 e(b_0,b_1)&=-\frac14,\\
 4e(b_j,x)&\equiv-\tau(x)\pmod4&&(x\in E).
 \end{aligned}
\end{equation}
\end{lemma}

\begin{proof}
The weighted rotation preserves each axis, hence its boundary homology
class.  Transfer of either axis boundary is a single boundary circle of
the upstairs plane-curve page.  In a positive plane-curve Seifert
convention, the pairing of two distinct complex branch boundaries is their
positive local intersection multiplicity.  The sum of all oriented
boundary classes is zero, so the self-pairing of the $V=0$ boundary is
$-(1+4a)=-n_0^2$, and that of $U=0$ is $-(1+4b)=-n_1^2$.
The two axes meet once.  Reversing the second orientation gives the
mixed value $-1$ upstairs.  Division by the covering degree proves the
first two formulas.

The integer-valued functional $x\mapsto4e(b_j,x)$ reduces modulo $4$
through $\tau$.  Indeed a loop whose $\tau$-class is zero lifts closed;
its full transfer consists of four deck translates, and the transferred
axis class is deck invariant.  The four Seifert pairings are equal, so
their sum is divisible by $4$.  A class in $\ker\tau$ can be represented
by such loop systems, and $\tau(b_j)=1$ generates $\mathbb Z/4$.
Evaluation at $b_j$ gives $4e(b_j,b_j)=-n_j^2\equiv-1\pmod4$,
proving the last congruence.
\end{proof}

\subsection{Two complementary local splittings}\label{subsec:arith-local}

We first make explicit why the Seifert operator of $e$ is defined.
Over $\mathbb Q$, transfer identifies $E_{\mathbb Q}$ with the
deck-invariant summand of the upstairs page homology. Let
$\widetilde s$ be the nondegenerate upstairs Milnor Seifert form
in the same ordered convention.
The deck transformations are $\widetilde s$-isometries, so their
averaging projection $p=\frac14\sum_{\gamma\in\mu_4}\gamma$ satisfies
$\widetilde s(px,y)=\widetilde s(x,py)$. Its image and kernel are
therefore orthogonal in both orders, and the restriction of
$\widetilde s$ to the invariant summand is nondegenerate. Under
transfer, $e$ is one fourth of this restriction,
so $e$ is nondegenerate and
$A_e=e^{-1}e^{\mathsf t}$ is the restriction of
$\widetilde A=\widetilde s^{-1}\widetilde s^{\mathsf t}$.
The weighted rotation $\widetilde M$ has rational weights $b/m,a/m$,
hence finite order, and commutes with the deck action. The identity
relating a plane-curve Seifert form and monodromy reads
$\widetilde A=\widetilde M^\varepsilon$ for one fixed
$\varepsilon\in\{1,-1\}$ in the common convention. The same
$\varepsilon$ applies to the base page. Thus $A_e$
has finite order and fixes the classes $b_j$, since the weighted
rotation fixes their transferred axis boundaries.

Set $v_j=b_j/n_j\in E_{\mathbb Q}$.  Since $A_eb_j=b_j$,
the form $e$ pairs $v_j$ symmetrically with every vector:
$e(v_j,x)=e(x,A_ev_j)=e(x,v_j)$.
Equation \eqref{eq:arith-boundary-pairings} gives
\begin{equation}\label{eq:arith-v-norm}
 e(v_j,v_j)=-\frac14.
\end{equation}
For a rational prime $\ell\nmid n_j$, write
$E_{\mathbb Z_\ell}=E\otimes\mathbb Z_\ell$ and let
\[
 \tau_{\mathbb Z_\ell}\colon E_{\mathbb Z_\ell}
 \longrightarrow(\mathbb Z/4)\otimes\mathbb Z_\ell
 \cong\mathbb Z_\ell/4\mathbb Z_\ell
\]
be the scalar extension of $\tau$. For $x\in E_{\mathbb Z_\ell}$,
the value $-4e(v_j,x)$ lies in $\mathbb Z_\ell$, since
$4e(b_j,x)$ is integral and $n_j$ is a unit. In this coefficient
ring, \eqref{eq:arith-boundary-pairings} gives
\[
 -4e(v_j,x)\equiv\tau_{\mathbb Z_\ell}(x)
       \pmod{4\mathbb Z_\ell}.
\]
The target is zero at odd primes. At $\ell=2$, the congruence uses
$n_j\equiv1\pmod4$ from Lemma~\ref{lem:simultaneous-parameters}.

\begin{lemma}[Local splitting]\label{lem:arith-local-splitting}
For every rational prime $\ell\nmid n_j$, the lattice
$E_{\mathbb Z_\ell}$ has the following $\mathbb Z_\ell$-module
splitting, orthogonal for the $\mathbb Q_\ell$-valued form $e$:
\[
 E_{\mathbb Z_\ell}=
 \mathbb Z_\ell v_j\perp U_{j,\ell},
 \qquad U_{j,\ell}\subset\ker(\tau_{\mathbb Z_\ell}).
\]
For $\nu=0,1$ this induces an integral Seifert isometry
\begin{equation}\label{eq:arith-local-alpha}
 \begin{split}
 \alpha_{\nu,j}\colon
 &(D_\nu,-s_{d_\nu})_{\mathbb Z_\ell}
       \perp(H\otimes U_{j,\ell},s_h\otimes e)
       \xrightarrow{\ \sim\ }(\Lambda_\nu,s_\nu)_{\mathbb Z_\ell},\\
 &(l,u)\longmapsto
 \bigl(\operatorname{pr}_\nu l,\ \pi_{\nu*}l\otimes v_j+u\bigr).
 \end{split}
\end{equation}
It also holds over $\mathbb Q$.
\end{lemma}

\begin{proof}
The coefficient of $x\in E_{\mathbb Z_\ell}$ on $v_j$ is
$c_j(x)=-4e(v_j,x)$.  It is integral because $4e(b_j,x)$ is integral and
$n_j$ is a unit in $\mathbb Z_\ell$, and it sends $v_j$ to $1$ by
\eqref{eq:arith-v-norm}.  Its kernel $U_{j,\ell}$ is therefore an
integral orthogonal summand. The local congruence above gives
$\tau_{\mathbb Z_\ell}(U_{j,\ell})=0$. At $\ell=2$,
$c_j(x)\equiv\tau_{\mathbb Z_2}(x)\pmod{4\mathbb Z_2}$
because $n_j\equiv1\pmod4$;
at odd primes the target of $\tau_{\mathbb Z_\ell}$ is zero.
Thus the projection $x\mapsto c_j(x)v_j$ is integral also at
$\ell=2$: all divisions are by the unit $n_j$.

The image of \eqref{eq:arith-local-alpha} satisfies the lattice
congruence \eqref{eq:arith-input-lattice}.  Conversely, write the second
coordinate of an arbitrary lattice element as
$z=x_H\otimes v_j+u$.  Its congruence says
$x_H-\pi_{\nu*}l\in4H_{\mathbb Z_\ell}$.  Replacing $l$ by
\[
 l+\pi_\nu^!\left((x_H-\pi_{\nu*}l)/4\right)
\]
is legitimate even at $\ell=2$, because the numerator belongs to
$4H_{\mathbb Z_\ell}$. It does not alter $\operatorname{pr}_\nu l$ and makes its image under
$\pi_{\nu*}$ equal to $x_H$.  This proves surjectivity; injectivity
follows from the rational direct sum.

The $H\otimes U_{j,\ell}$ summand is orthogonal to the image of $D_\nu$.
On two vectors from $D_\nu$, transfer adjunction gives
\[
 s_{d_\nu}(\operatorname{pr}_\nu l,
                  \operatorname{pr}_\nu l')
 =s_{d_\nu}(l,l')-
        \frac14s_h(\pi_{\nu*}l,\pi_{\nu*}l').
\]
The product-form contribution on $v_j$ is the negative of the last
fraction, by $e(v_j,v_j)=-1/4$.  Combined with the leading minus sign
on $K_\nu$ in \eqref{eq:arith-input-lattice}, the result is
$-s_{d_\nu}(l,l')$, as claimed.
\end{proof}

Use the fixed plane isometry $\mathcal I$ from
\eqref{eq:arith-plane-isometry} on the $D$-summands and the identity
on $H\otimes U_{j,\mathbb Q}$.  The rational versions of
\eqref{eq:arith-local-alpha} give Seifert isometries
\begin{equation}\label{eq:arith-two-J}
 J_j=\alpha_{1,j}(\mathcal I\oplus1)\alpha_{0,j}^{-1}
 \colon(F_0,s_0)\xrightarrow{\ \sim\ }(F_1,s_1).
\end{equation}
They are lattice isometries at every $\ell\nmid n_j$.  No compatibility
between $\mathcal I$ and the two maps $\pi_\nu$ is required: those maps enter
only through the explicit changes of coordinates $\alpha_{\nu,j}$.
Since $\gcd(n_0,n_1)=1$, the two isometries $J_0,J_1$
cover all primes locally.

\subsection{The discrepancy on two primary sectors}
\label{subsec:arith-correction}

The local choices will give one integral isometry if their discrepancy
\[
 \Delta=J_0^{-1}J_1\colon F_0\longrightarrow F_0
\]
can be corrected at the primes dividing $n_0$ without disturbing the
source lattice at the other primes. We first identify the rational
isometry groups containing this discrepancy. The normalized boundary
vectors $v_0,v_1$ are exchanged by a reflection; comparison of the corresponding
reflections on $F_0$ and $F_1$ will remove the determinant and spinor
obstructions. To state that comparison, we describe the two monodromy
sectors on which the reflections act.

Let $A_\nu=s_\nu^{-1}s_\nu^{\mathsf t}$ be the Seifert operator.
It is an isometry of $s_\nu$.  By the rational orthogonal decomposition
\eqref{eq:arith-input-lattice}, its action is the operator of $d_\nu$
on $K_\nu$ and the tensor product of the operators of $h$ and $e$
on $H\otimes E$.  The operator of $e$ has finite order because
\eqref{eq:arith-upstairs-page} is weighted homogeneous; the operator
of $h$ has order $5$.  The plane-curve semisimplicity input therefore
shows that $A_\nu$ is semisimple at $1$ and at primitive fifth roots.
Let
\[
 F_{\nu,1}=\ker(A_\nu-1),\qquad
 F_{\nu,5}=\ker\Phi_5(A_\nu),\qquad
 B_\nu=s_\nu+s_\nu^{\mathsf t}.
\]
These primary spaces split off rationally and orthogonally, because
their eigenvalue sets are stable under inversion.  On them $B_\nu$
is nondegenerate: $s_\nu^{\mathsf t}=s_\nu A_\nu$, hence
$B_\nu=s_\nu(1+A_\nu)$, and neither primary sector has eigenvalue
$-1$.  On $F_{\nu,1}$ it is a rational quadratic form.  On
$F_{\nu,5}$ the action of $A_\nu$ makes the space a vector space over
$L=\mathbb Q(\zeta)$, $\zeta=e^{2\pi i/5}$, with
$L^+=\mathbb Q(\zeta+\zeta^{-1})$. Define the $L$-valued form
$\mathfrak h_{\nu,5}$ on $F_{\nu,5}$ by trace duality:
\[
 \operatorname{Tr}_{L/\mathbb Q}\!\left(c\,\mathfrak h_{\nu,5}(x,y)\right)
   =B_\nu(cx,y)\qquad(c\in L).
\]
The $B_\nu$-adjoint of $c$ is $\bar c$, because $A_\nu$ is a
$B_\nu$-isometry. Nondegeneracy of the field trace pairing therefore
gives sesquilinearity and Hermitian symmetry for $L/L^+$, and
nondegeneracy follows from that of $B_\nu$. The same identity shows
that every $L$-linear $B_\nu$-isometry preserves $\mathfrak h_{\nu,5}$.

For the base form on $H_{\mathbb Q}$, set
\[
 A_h=s_h^{-1}s_h^{\mathsf t},\qquad
 H^{(1)}=\ker(A_h-1),\qquad H^{(5)}=\ker\Phi_5(A_h).
\]
With the same convention as above, $A_h=M_h^\varepsilon$.

\begin{lemma}[The discrepancy lifts]\label{lem:arith-discrepancy}
The rational Seifert isometry
$\Delta=J_0^{-1}J_1$ acts as the identity outside
$F_{0,1}\oplus F_{0,5}$.  Its restriction to $F_{0,1}$ lifts to the
rational points of the spin group of $B_0|F_{0,1}$.  Its restriction to
$F_{0,5}$ lies in the special unitary group of the $L/L^+$ Hermitian
form.
\end{lemma}

\begin{proof}
The two $v_j$ have the same norm $-1/4$, while
$e(v_0,v_1)=-1/(4n_0n_1)$.  Therefore
\[
 w=v_0-v_1,\qquad
 e(w,w)=-\frac12\left(1-\frac1{n_0n_1}\right)\ne0.
\]
The vector $w$ is monodromy fixed and hence pairs symmetrically
for $e$.  Reflection
\[
 R_w(x)=x-\frac{2e(x,w)}{e(w,w)}w
\]
is consequently an $e$-isometry commuting with monodromy, and it
exchanges $v_0,v_1$.  Define $\rho_\nu$ on $F_\nu$ as identity on
$(K_\nu)_{\mathbb Q}$ and $1_H\otimes R_w$ on
$(H\otimes E)_{\mathbb Q}$.  This is a rational Seifert
involution.  Let $\mathcal R_\nu$ be identity on $D_\nu$ and
$1_H\otimes R_w$ from $H\otimes U_{0,\mathbb Q}$ to
$H\otimes U_{1,\mathbb Q}$.  Equation
\eqref{eq:arith-local-alpha} gives
$\rho_\nu\alpha_{\nu,0}=\alpha_{\nu,1}\mathcal R_\nu$.
The same $R_w$ occurs for both coverings, so the square formed by
$\mathcal R_0,\mathcal R_1$ and $\mathcal I\oplus1$ commutes.  Consequently
\begin{equation}\label{eq:arith-J-conjugacy}
 J_1=\rho_1J_0\rho_0^{-1},\qquad
 \Delta=(J_0^{-1}\rho_1J_0)\rho_0^{-1}.
\end{equation}
The minus space of $\rho_\nu$ is
$W_\nu=H_{\mathbb Q}\otimes\mathbb Qw$. Its restricted Seifert form
is $e(w,w)s_h$, so its Seifert operator is $A_h$. In particular it
lies entirely in $F_{\nu,1}\oplus F_{\nu,5}$. Both factors of
\eqref{eq:arith-J-conjugacy}, and therefore $\Delta$, act
identically on every other rational primary summand.

The minus spaces within the two primary sectors are
\[
 W_{\nu,1}=H^{(1)}\otimes\mathbb Qw,\qquad
 W_{\nu,5}=H^{(5)}\otimes\mathbb Qw.
\]
The canonical map $x\otimes w\mapsto x\otimes w$ from $W_0$ to
$W_1$ preserves the restricted Seifert and symmetric forms and
intertwines $A_h$. It therefore restricts to isometries on the
individual sectors displayed above. These restrictions are
nondegenerate: $e(w,w)\ne0$, and $s_h+s_h^{\mathsf t}$ is
nondegenerate on both $H^{(1)}$ and $H^{(5)}$. Since the Seifert
isometry $J_0$ intertwines $A_0$ and $A_1$, composition with
$J_0^{-1}$ identifies, separately on each sector, the minus spaces
of $\rho_0$ and $J_0^{-1}\rho_1J_0$.
On $F_{0,5}$ the isometry $\Delta$ commutes with the $L$-action
and preserves the Hermitian form recovered from $B_0$.  The two
involutions in \eqref{eq:arith-J-conjugacy} have the same
$L$-dimension of their minus spaces, so their $L$-determinants
cancel: $\det_L\Delta=1$.
It is therefore an $L^+$-rational special-unitary element.

On $F_{0,1}$ the same comparison gives determinant $1$.  The
spinor norm of an involution with nondegenerate minus space is the
product, modulo rational squares, of the quadratic values of an
orthogonal basis of that space.  The minus spaces just identified
are isometric, so their spinor norms are
equal and cancel in the product.  Thus $\Delta$ has trivial
spinor norm.  The standard exact sequence for
$\operatorname{Spin}\to\operatorname{SO}$ now supplies a rational
spin lift
\cite[Section~C.5, paragraph before Example~C.5.5 and Lemma~C.5.8]{Conrad2014Reductive}.
\end{proof}

\subsection{Noncompact real factors}\label{subsec:arith-real}

Lemma~\ref{lem:arith-discrepancy} places the required correction in a
spin group over $\mathbb Q$ and a special unitary group over $L^+$.
Strong approximation in these groups requires noncompactness at an
archimedean place. We now prove the corresponding indefiniteness of
the quadratic and Hermitian forms. Both signs come from holomorphic
and antiholomorphic classes of the cone page.

\begin{lemma}[Real signs]\label{lem:arith-real-signs}
The quadratic form $B_\nu|F_{\nu,1}$ is indefinite and has dimension at
least five.  The Hermitian form on $F_{\nu,5}$ has dimension at least two
over $L$ and is indefinite at the real place of $L^+$ indexed by
the pair $\{\zeta,\bar\zeta\}$. Both assertions hold for $\nu=0,1$.
\end{lemma}

\begin{proof}
By Lemma~\ref{lem:plane-data}, the rational primary parts
$H^{(1)}$ and $H^{(5)}$ have dimensions two and four, respectively,
and together equal
$H_{\mathbb Q}$. Their eigenvalue sets are stable under inversion, so
they are nondegenerate orthogonal summands for $s_h$. The restriction
to $H^{(1)}$ is symmetric, and each complex primitive fifth
eigenspace is one-dimensional and pairs nondegenerately with its
inverse eigenspace.

To obtain both Hodge signs in the primitive fifth eigenspaces of
$E_{\mathbb C}$, we construct holomorphic forms with inverse primitive
fifth characters. Their conjugate antiholomorphic classes will give
opposite cup-product signs. Set $m=4ab+a+b$, and write
$a=ga'$, $b=gb'$ with $\gcd(a',b')=1$.  The weights of $U,V$ in
\eqref{eq:arith-upstairs-page} are $b/m,a/m$.  For $q=1$ and $q=4$
there are positive integers $i,j$, congruent modulo $4$, with
\begin{equation}\label{eq:arith-weight-indices}
 \frac{bi+aj}{m}=\frac q5.
\end{equation}
The existence and positivity of these two index pairs were proved in
Lemma~\ref{lem:parameter-indices}.

On the smooth affine curve $\widetilde P=1$, the differential
\begin{equation}\label{eq:arith-eigenform}
 \omega_{i,j}=U^{i-1}V^{j-1}
        \frac{dU}{\partial_V\widetilde P}
       =-U^{i-1}V^{j-1}
        \frac{dV}{\partial_U\widetilde P}
\end{equation}
is regular and nonzero.  The second expression is used where
$\partial_V\widetilde P=0$.  Its deck character is
$\omega^{i-j}=1$, so it descends to the quotient, and its weighted
rotation character is $e^{2\pi iq/5}$.
To obtain the needed holomorphic classes, we check that these
differentials extend across every end of the normalized
compactification. At an end put
$u=\operatorname{ord}U$, $v=\operatorname{ord}V$.
If $4au<4bv$, then $u<0$ and $v=-(4a+1)u$; the order of
\eqref{eq:arith-eigenform} is
$(i-(4a+1)j)u-1\ge0$.  Here
$i<4a+1$, because \eqref{eq:arith-weight-indices}, $q\le4$ and
$j\ge1$ give $bi\le4m/5-a<b(4a+1)$.
The case $4bv<4au$ is symmetric, using the second expression and
$j<4b+1$.  In the remaining case $4au=4bv<0$, leading terms in
$U^{4a}+V^{4b}$ cancel and
$\operatorname{ord}(\partial_V\widetilde P)=u+4bv$.
There is a positive integer $c$ with $(u,v)=-c(b',a')$; the order is
\[
 (i-1)u+(j-1-4b)v-1
 =c\left(1-\frac q5\right)\frac mg-1\ge0.
\]
These are all possible ends, since one of $U,V$ has a pole.

The two differentials for $q=1,4$ thus give nonzero holomorphic
forms on the smooth compactification with characters $\zeta$ and
$\bar\zeta$.  Complex conjugation gives antiholomorphic forms of
the opposite characters.  The cup-product Hermitian form has
opposite signs on holomorphic and antiholomorphic forms. Poincar\'e
duality exchanges reciprocal characters and gives opposite signs for the
Hermitian intersection form on the corresponding homology classes. Since both characters
have been constructed, each of the two homology eigenspaces has both
signs. With complex coefficients, the map from affine-page homology to
compactification homology is surjective and preserves intersections. Average a linear section
of this map over the finite group generated by the commuting deck
transformations and weighted rotation. The resulting equivariant
section lifts the classes with their characters and intersection
numbers unchanged. Transfer over $\mathbb C$ identifies
the deck-invariant part with $E_{\mathbb C}$ and preserves the
monodromy characters.  Since $e-e^{\mathsf t}$ is, up to our common
Seifert convention, the page intersection form, the Hermitian form
$i(e-e^{\mathsf t})(x,\bar x)$ has both signs in each of
$E_\zeta$ and $E_{\bar\zeta}$.

We transfer these signs to the fixed quadratic sector by tensoring
a primitive base eigenspace with the inverse eigenspace of $E$.
After complexifying, choose $\lambda\in\{\zeta,\bar\zeta\}$ and a
nonzero $x\in\ker(A_h-\lambda)\subset H_{\mathbb C}$. This eigenspace is one-dimensional and
pairs nondegenerately with its inverse eigenspace. The preceding
statement that both signs occur applies to
$\ker(A_e-\lambda^{-1})\subset E_{\mathbb C}$
under either inverse monodromy convention. For $y$ in this eigenspace, put
\[
 I_H=i(s_h-s_h^{\mathsf t})(x,\bar x),\qquad
 I_E=i(e-e^{\mathsf t})(y,\bar y).
\]
These are real, and $I_H\ne0$. The relation $s^{\mathsf t}=sA$ in
each factor, together with the vanishing of $s_h(x,x)$ and $e(y,y)$
when $\lambda^2\ne1$, gives
\[
 B_\nu(x\otimes y+\bar x\otimes\bar y,\,
       x\otimes y+\bar x\otimes\bar y)
 =-\frac{4I_HI_E}{|1-\lambda|^2}.
\]
The vectors in this display lie in $F_{\nu,1}\otimes_{\mathbb Q}\mathbb R$, and the two
signs of $I_E$ prove that $B_\nu|F_{\nu,1}$ is indefinite. These
tensors and their conjugates already contribute at least four real dimensions.
The invariant part of $H$ has dimension two, and $b_0$ is a nonzero
invariant vector of $E$ by \eqref{eq:arith-boundary-pairings}; its
tensor contributes at least two more dimensions.

For the Hermitian sector, we tensor the primitive eigenspace of $E$
with a fixed base vector. Choose $x\in H\cap H^{(1)}$ with $s_h(x,x)\ne0$.
Such an integral vector exists: the nondegenerate symmetric
restriction to $H^{(1)}$ has a rational vector of nonzero norm, and
clearing denominators puts it in $H$. Tensor $x$ with $E_\zeta$.
This lies in the
$\zeta$ component of $F_{\nu,5}$ for one of the two inverse
monodromy conventions.  If the eigenvalue of $A_e$ on $E_\zeta$
is $\lambda\in\{\zeta,\zeta^{-1}\}$, then on this eigenspace
\[
 (e+e^{\mathsf t})(y,\bar y)
 =\frac{1+\lambda^{-1}}{i(1-\lambda^{-1})}
       i(e-e^{\mathsf t})(y,\bar y),
\]
and the displayed ratio is a nonzero real number.  The two signs
found above therefore give both signs for the Hermitian form of
$B_\nu$ at this real embedding.  Its dimension is at least two.
\end{proof}

\subsection{A single integral isometry}\label{subsec:arith-approximation}

We can now combine the rational lift of the discrepancy with the real
noncompactness just proved. The approximation will be imposed inside
stabilizers of the full lattice $\Lambda_0$, so it preserves the
mod-four gluing even though the rational primary decomposition need
not split that lattice integrally.

\begin{proposition}[Integral Seifert congruence]
\label{prop:arith-integral-congruence}
For the two germs constructed above, the two integral Milnor Seifert
forms $(\Lambda_0,s_0)$ and $(\Lambda_1,s_1)$ are isometric.
\end{proposition}

\begin{proof}
By Lemma~\ref{lem:arith-real-signs},
$G_1=\operatorname{Spin}(F_{0,1},B_0)$ is a connected, simply
connected, absolutely almost simple group over $\mathbb Q$ of
quadratic dimension at least five
\cite[Proposition~C.3.10, Lemma~C.4.1, Proposition~C.4.10, and Example~C.6.6]{Conrad2014Reductive}.
The proved indefiniteness makes it noncompact over $\mathbb R$.
The second correction group
$G_5=\operatorname{SU}(F_{0,5},\mathfrak h_{0,5})$
is connected, simply connected,
and absolutely almost simple over $L^+$: after algebraic closure it
is $\operatorname{SL}_d$, with $d\ge2$
\cite[Proposition~24.40(b), equation~(158), and Summary~21.96]{Milne2022}.
It is noncompact at the real place found in
Lemma~\ref{lem:arith-real-signs}.
We use
$G=G_1\times\operatorname{Res}_{L^+/\mathbb Q}G_5$,
acting on the indicated primary spaces and as identity on their
$s_0$-orthogonal primary complement. The Spin action preserves
$B_0$ and commutes with $A_0=1$ on $F_{0,1}$. The special-unitary
action preserves the Hermitian form, hence $B_0$, and commutes with
$A_0$ on $F_{0,5}$ by $L$-linearity. On both sectors $1+A_0$ is
invertible and
\[
 s_0=B_0(1+A_0)^{-1}.
\]
Consequently this action of $G$ preserves the Seifert form $s_0$.
We obtain the density assertion for $G$ by treating its two factors
separately.

The strong-approximation theorem of Kneser and Platonov, in the
precise form of \cite[Theorem~2.3]{Rapinchuk2014}, says that for a
connected, simply connected, absolutely almost simple group over a
number field, noncompactness at a place in the omitted finite set
implies density of rational points in the adeles away from that set.
Apply it to $G_1$ over $\mathbb Q$ and to $G_5$ over $L^+$, omitting
all archimedean places.  Restriction of scalars and a finite
product give density of $G(\mathbb Q)$ in $G(\mathbb A_f)$.

Write $\rho\colon G\to\operatorname{GL}(F_0)$ for this rational
action.  Lemma~\ref{lem:arith-discrepancy} gives a rational lift
$\widetilde\Delta\in G(\mathbb Q)$ with
$\rho(\widetilde\Delta)=\Delta$.  For each rational prime $\ell$,
put $a_\ell=\widetilde\Delta$ if $\ell\mid n_0$, and
$a_\ell=1$ otherwise. The tuple $\mathbf a=(a_\ell)_\ell$ has finite
support and lies in $G(\mathbb A_f)$. We seek one rational point
$g_{\mathrm{rat}}\in G(\mathbb Q)$ whose localizations
$g_{\mathrm{rat},\ell}$ differ from these selected corrections by
automorphisms of the full source lattice:
\[
 \rho(a_\ell^{-1}g_{\mathrm{rat},\ell})
       (\Lambda_0\otimes\mathbb Z_\ell)
   =\Lambda_0\otimes\mathbb Z_\ell
       \quad\text{for every }\ell.
\]
To construct a suitable open neighborhood of $\mathbf a$, define the
exact source-lattice stabilizer
\[
 C_\ell=\left\{\gamma\in G(\mathbb Q_\ell):
   \rho(\gamma)(\Lambda_0\otimes\mathbb Z_\ell)
          =\Lambda_0\otimes\mathbb Z_\ell\right\}.
\]
Each $C_\ell$ is compact open.  Here the spin action has finite
central kernel, so it need not be faithful; compactness follows
from finiteness of the covering onto its algebraic image.
For the adelic topology we only need compact opens inside these
stabilizers.  Spread out $G$ and its rational action on $F_0$ to a
smooth affine group scheme $\mathcal G$ over $\mathbb Z[1/N]$,
enlarging $N$ so that the action preserves
$\Lambda_0\otimes\mathbb Z[1/N]$.  At every $\ell\nmid N$ its
standard compact open $\mathcal G(\mathbb Z_\ell)$ lies in $C_\ell$.
For each of the finitely many primes $\ell\mid N$, choose a compact
open subgroup $K_\ell\subset C_\ell$; at the other primes put
$K_\ell=\mathcal G(\mathbb Z_\ell)$.  Then
$\mathcal K=\prod_\ell K_\ell$ is compact open in $G(\mathbb A_f)$ and
$\mathbf a\mathcal K$ is an open coset contained in the desired set of local
integral corrections.

At $\ell\nmid n_0$, Lemma~\ref{lem:arith-local-splitting}
makes $J_0$ an integral lattice isometry.  At $\ell\mid n_0$
we have $\ell\nmid n_1$, so $J_0\rho(a_\ell)=J_1$ is integral.
Consequently
\[
 J_0\rho(a_\ell)(\Lambda_0\otimes\mathbb Z_\ell)
                  =\Lambda_1\otimes\mathbb Z_\ell
       \quad\text{for every }\ell.
\]
Density supplies $g_{\mathrm{rat}}\in G(\mathbb Q)\cap\mathbf a\mathcal K$.
Writing $g_{\mathrm{rat},\ell}=a_\ell k_\ell$
with $k_\ell\in K_\ell\subset C_\ell$
shows that
$J_0\rho(g_{\mathrm{rat}})$ induces the same local lattice equality at every prime.
A full $\mathbb Z$-lattice in a rational vector space is the
intersection of its localizations at all primes; applying this to
both sides yields
$J_0\rho(g_{\mathrm{rat}})\Lambda_0=\Lambda_1$.
Since $J_0$ and $\rho(g_{\mathrm{rat}})$ preserve their respective Seifert forms,
this is the asserted integral congruence.
\end{proof}

\subsection{The embedded links}\label{subsec:arith-links}

The final step uses Kato's Theorem~B and his definition in Section~2
\cite[pp.~454--455]{Kato1974}. The theorem
has no homotopy-sphere hypothesis on the binding. In his terminology,
a spinnable structure
on an oriented $(2n+1)$-manifold is an open book with compact
$2n$-dimensional page $F$, characteristic diffeomorphism fixed
near $\partial F$, and an identification of its open-book quotient
with the ambient manifold.  It is simple when $F$ is obtained from
a ball by handles of index at most $n$.  Theorem~B states that for
$n\ge3$ two such structures on $S^{2n+1}$ with congruent integral
Seifert matrices are isomorphic.  Kato explicitly notes on p.~454
that this applies to isolated complex hypersurface singularities;
his footnote credits Durfee's independent result
\cite{Durfee1974}.

For each function constructed above, the origin is an isolated
critical point in a smooth complex four-dimensional chart.  Its
Milnor open book lies on the standard $S^7$, so $n=3$; its compact
page has real dimension six and its binding has real dimension five.
We use the analytic form of Milnor's theorem, recalled explicitly by
N\'emethi \cite[pp.~63--64]{Nemethi1999SeifertInfinity}: an isolated
analytic germ supplies a simple spinnable Milnor structure. In our
four-dimensional charts this gives the handle bound required by Kato.

We also specify the characteristic diffeomorphism near the binding.
The Milnor open book has the local product model
$\partial F\times D^2$, with the fibration given by the angular
coordinate on $D^2\setminus\{0\}$. Choose the horizontal angular
field to be the product rotation on a smaller such neighborhood,
and extend it over the complement by a partition of unity among
horizontal lifts. Its return after one turn fixes a collar of
$\partial F$ pointwise.

Theorem~\ref{thm:tube} identifies \eqref{eq:arith-input-lattice}
with the actual integral Milnor Seifert forms in the same ordered
push-off convention. Proposition~\ref{prop:arith-integral-congruence}
is an onto isometry of these integral lattices. In integral bases
of the page groups $H_3(F;\mathbb Z)$, its matrix is therefore
unimodular and gives exactly an integral congruence of Seifert
matrices. The common complex and positive-target orientations
make any universal sign or transpose in the conversion to Kato's
convention the same on both sides, so the congruence is preserved.
Kato's Theorem~B now gives an orientation-preserving diffeomorphism
$\phi$ of the two sphere-link pairs: the isomorphism of open books
carries pages and therefore their embedded boundaries.

To pass to the germs, use the analytic ambient-pair cone theorem,
attributed to \L{}ojasiewicz in
\cite[Theorem~1.4]{Massey2015LeCycles}. It identifies each sufficiently
small ball pair, including its center, with the cone on its
sphere-link pair. The map
$[u,t]\mapsto[\phi(u),t]$ is a homeomorphism of the two cone pairs.
Conjugating it by the two analytic pair-cone homeomorphisms gives
the required homeomorphism germ of ambient pairs.

\begin{proof}[Proof of Theorem~\ref{thm:main}]
Proposition~\ref{prop:ambient-orders} and Lemma~\ref{lem:isolation}
give smooth ambient fourfold germs and reduced isolated functions of
orders five and four. Theorem~\ref{thm:tube} identifies their actual
integral Milnor Seifert forms, and Proposition~\ref{prop:arith-integral-congruence}
gives an integral congruence. The Kato argument above identifies their
embedded links. The analytic pair-cone construction above extends that equivalence
to an ambient homeomorphism germ, proving the assertion.
\end{proof}
